\section*{2. 有限定域上的余层次}

\subsection*{2.0.}
设 $X$ 是 $\mathbb F_q$ 上的有限型分离概形. 对每个素数
$\ell\ne\operatorname{char}(\mathbb F_q)$, 记
\begin{equation}\tag{2.0.1}
H_c^i(\overline X,\mathbb Q_\ell)
\end{equation}
为 $\overline X=X\otimes_{\mathbb F_q}\overline{\mathbb F}_q$ 的紧支撑上同调群.
它们是带有 Galois 群 $\operatorname{Gal}(\overline{\mathbb F}_q/\mathbb F_q)$ 作用的
有限维 $\mathbb Q_\ell$ 向量空间. 记
\begin{equation}\tag{2.0.2}
\varphi_q\in\operatorname{Gal}(\overline{\mathbb F}_q/\mathbb F_q)
\end{equation}
为 Frobenius 自同构. 已知 (SGA~5 XIII~3.2.1), $X$ 的 zeta 函数 (XXII~3.0.1) 由下式给出
\begin{equation}\tag{2.0.3}
Z(t,X/\mathbb F_q)=
\prod_{i\geq0}
\det\!\left(1-t\varphi_q^{-1}\mid
H_c^i(\overline X,\mathbb Q_\ell)\right)^{(-1)^{i+1}}.
\end{equation}
Deligne 最近推广了 Washnitzer 对光滑射影 $X$ 的结果, 证明了下列定理;
其证明将在后面的附录 (5.5.3) 中给出.

% French-source verification: original PDF page 387, printed folio 377,
% Expose XXI local page 14, has q^{i-\dim X}; the loose English comma is a transcription error.
\subsubsection*{定理 2.1.}
$\varphi_q^{-1}$ 在 $H_c^i(\overline X,\mathbb Q_\ell)$ 上的特征值都是代数整数;
若 $i>\dim X$, 则它们作为代数整数均被 $q^{i-\dim X}$ 整除.

特别地, 直接按定义得到:

\medskip
\noindent\textbf{(2.1 bis).}
$\varphi_q^{-1}$ 在 $H_c^i(\overline X,\mathbb Q_\ell(-j))$ 上的特征值
均为被 $q^j$ 整除的代数整数.

\subsubsection*{推论 2.2.}
设 $X$ 在 $\mathbb F_q$ 上固有、光滑且几何连通. 则 $\varphi_q^{-1}$ 在
$N^jH^q(\overline X,\mathbb Q_\ell)$ (记号见 XX~2.0.6) 上的特征值
均为被 $q^j$ 整除的代数整数.

\medskip
\noindent\emph{证明.}
此断言在有限扩张 $\mathbb F_{q^M}/\mathbb F_q$ 下不变; 这种扩张相当于将所考察的
特征值替换为其 $M$ 次幂. 作此基变换后, 由定义可得有限多个 $\mathbb F_q$ 上的
射影、光滑、几何连通概形 $T_i$, 其维数为 $d_i$, 以及
$T_i\times_{\mathbb F_q}X$ 上余维为 $j+d_i$ 的代数圈 $Z_i$, 使得
$\operatorname{Gal}(\overline{\mathbb F}_q/\mathbb F_q)$ 模同态
\begin{equation*}
\bigoplus_i H^{q-2j}(\overline T_i,\mathbb Q_\ell(-j))
\xrightarrow{\ \oplus_i\operatorname{Cl}(Z_i)\ }
N^jH^q(\overline X,\mathbb Q_\ell)
\end{equation*}
为满射. 结论由 (2.1 bis) 得出.

\medskip
\noindent\emph{底本注.}
照录法文层保留了印本措辞 \emph{$\mathbb Q_\ell$-vectoriels}、
\emph{fonction zêta ... donné}、扩张 $\mathbb F_{q^M}/\mathbb F_q$ 之后被涂改的短语、
分子留空的 Galois 群以及未明确写直和号的定义域.
校订法文层和英文层采用标准的向量空间术语及语法一致形式, 明确写出取 $M$ 次幂的运算,
在 Galois 群中补回 $\overline{\mathbb F}_q$, 并写出 $\oplus_i\operatorname{Cl}(Z_i)$
所要求的直和.
